\section{Apps, Modules, and the Store}
\label{sec:apps-modules}

\subsection{App Ecosystem}
\label{sec:app-ecosystem}

The application layer divides cleanly into two categories: first-party
applications built as part of this project, and the broader ecosystem of
third-party Linux software. Both operate under the same permission system and
benefit from the same infrastructure, but they differ in how deeply they
integrate with the platform.

\paragraph{First-party applications.}
All core system applications are built with the standard stack: Tauri, ui-kit,
and os-sdk. Each one follows the shared design language, implements an MCP
server for AI integration, emits structured events to the Event Bus, and uses
the shared TOML configuration system. Table~\ref{tab:first-party} lists the
planned core applications.

\begin{table}[htbp]
\caption{First-party system applications}
\label{tab:first-party}
\begin{tabularx}{\linewidth}{lX}
\toprule
\textbf{Application} & \textbf{Notes} \\
\midrule
File Manager    & Graph-aware features: recent files by project, related files sidebar \\
Terminal        & MCP server exposing command history and session output \\
Settings        & System-wide configuration: theming, permissions, AI level, accounts \\
Store           & App and module discovery, installation, and updates \\
Text Editor     & Basic editor primarily for configuration files \\
System Monitor  & Process list, resource usage, and Anomaly Detector alerts \\
Wine Manager    & Prefix management and \emph{Winetricks}~\cite{winetricks} access. A Wine \emph{prefix} is an isolated directory tree containing a synthetic Windows file system and registry; separate prefixes let different applications use different Wine versions or installed Windows libraries. Winetricks is a helper script that automates installation of common Windows runtimes (DirectX, Visual C++ Redistributables, .NET) and popular applications into a prefix. \\
\bottomrule
\end{tabularx}
\end{table}

\paragraph{Third-party integration layers.}
Third-party applications are not rewritten. Integration happens in four layers
of increasing depth, each building on the one below.

At the baseline, every application on the system receives passive tracking via
eBPF regardless of whether it does anything to support this platform: file
access, network connections, and process lifecycle are recorded automatically.
No changes to the application are required.

For applications that produce useful log output, a wrapper captures stdout and
stderr and passes them through asynchronous LLM extraction to produce graph
entities. These wrappers are shipped alongside packages in the package manager
where practical, and the extraction runs in the background without blocking
anything.

For popular applications, dedicated MCP server wrappers expose structured
interfaces to the AI layer. These are maintained as separate packages and
where possible contributed back to the community as independently useful tools.

Finally, applications that explicitly support the platform APIs emit structured
typed events directly to the Event Bus and implement native MCP servers. This
is the long-term goal for major open-source applications but is not a
precondition for any of the layers above it.

\paragraph{Shared standards.}
Applications that choose to integrate with the platform work with a common set
of conventions: TOML configuration files under
\texttt{\textasciitilde/.config/<app-id>/}, CSS custom properties from the
token system for applications that render via WebView, the standard
notification daemon via the Freedesktop D-Bus interface or the Event Bus for
richer notifications, and a published MCP server interface specification.

\subsection{Notification System}
\label{sec:notifications}

The Notification Daemon is the single entry point for all notifications from
all sources. Figure~\ref{fig:notifications} shows the pipeline.

\begin{figure}[htbp]
\centering
\begin{tikzpicture}[
  src/.style={
    draw, thick, rectangle, rounded corners=2pt,
    minimum width=3.0cm, minimum height=0.8cm,
    font=\footnotesize, align=center, inner sep=4pt
  },
  proc/.style={
    draw, thick, rectangle,
    minimum width=3.0cm, minimum height=0.8cm,
    font=\footnotesize, align=center, inner sep=4pt
  },
  store/.style={
    draw, thick, cylinder, shape border rotate=90,
    minimum width=1.8cm, minimum height=0.8cm,
    font=\footnotesize, align=center, inner sep=4pt
  },
  arr/.style={->, thick},
]

\node[src]  (dbus)   at (0,  2.4) {libnotify\\(D-Bus standard)};
\node[src]  (native) at (0,  1.2) {Native apps\\(Event Bus)};
\node[src]  (system) at (0,  0.0) {System components\\(Anomaly Detector\ldots)};

\node[proc] (nd)     at (3.8, 1.2) {Notification\\Daemon};

\node[proc]  (toast) at (7.4,  2.4) {Toast Renderer\\(Tauri)};
\node[proc]  (nc)    at (7.4,  1.2) {Notification Center\\(Tauri shell)};
\node[store] (al)    at (7.4,  0.0) {Audit Log};

\draw[arr] (dbus)   -- (nd);
\draw[arr] (native) -- (nd);
\draw[arr] (system) -- (nd);
\draw[arr] (nd)     -- (toast);
\draw[arr] (nd)     -- (nc);
\draw[arr] (nd)     -- (al);

\end{tikzpicture}
\caption{Notification pipeline. All sources converge on the Notification Daemon,
which normalizes, groups, and applies priority logic before forwarding to the
toast renderer and notification center. Every notification is logged in the
Audit Log by sender and timestamp, never by content.}
\label{fig:notifications}
\end{figure}

The daemon receives from two sources.

The first is the standard \emph{Freedesktop Notifications}
interface~\cite{fdo:notifications}: a D-Bus service registered at
\texttt{org.freedesktop.Notifications} that any application can send a
notification to by calling \texttt{Notify} with a summary, body, and urgency
level. \emph{libnotify}~\cite{libnotify} is the C library that wraps this
D-Bus call; \texttt{notify-send} is its command-line frontend. Because this
interface has been the Linux notification standard since 2006, every third-party
application that shows notifications already uses it and requires no changes.

The second source is the Event Bus, through which native system components send
richer notifications with priority metadata and graph context. Both sources are
normalized to a common internal schema before any further processing.

\begin{table}[htbp]
\caption{Notification priority levels}
\label{tab:notifications}
\begin{tabularx}{\linewidth}{lXX}
\toprule
\textbf{Level} & \textbf{Behavior} & \textbf{Do Not Disturb} \\
\midrule
Critical & Persistent toast, no auto-close &
  Always shown, never silenced \\
High     & Toast auto-closes after approximately 8 seconds &
  Delayed, not blocked \\
Normal   & Toast auto-closes after approximately 4 seconds &
  Silenced \\
Low      & No toast; goes silently to Notification Center only &
  Silenced \\
\bottomrule
\end{tabularx}
\end{table}

Critical priority is reserved for system-level alerts only: Anomaly Detector
warnings, component compatibility failures, and security events. Applications
cannot self-assign Critical. Notifications from the same application are
grouped in the Notification Center; as a toast, only the most recent
notification from an application is shown, with a counter if a previous one
from the same application is still visible. Do Not Disturb has three modes
(off, on, and scheduled by time of day) and is togglable from the taskbar quick
panel.

\subsection{Permission System}
\label{sec:permissions}

Every application on the system, whether a first-party app, a Flatpak, or an
RPM package, operates within a permission profile that defines what it can
access. The system is unified: one file format, one enforcement layer, one
place for the user to review and adjust everything.

\paragraph{Permission file format.}
Each application has a TOML permission file at
\texttt{\textasciitilde/.config/permissions/<app-id>.toml} that serves as the
sole source of truth for what that application is allowed to do. The file is
human-readable and user-editable; the Settings app provides a UI for those who
prefer not to edit it directly.

Permissions are enforced through three complementary Linux kernel
mechanisms~\cite{archwiki:namespaces}. \emph{Linux namespaces} partition kernel
resources: a process in its own mount namespace sees only the filesystem
subtrees explicitly provided to it; a process in its own network namespace sees
only the network interfaces and connections assigned to it. \emph{AppArmor}
is a Linux Security Module that enforces mandatory access control (MAC) via
profiles: the kernel checks every file open, network connect, and capability
use against the AppArmor profile and blocks anything not listed~\cite{archwiki:apparmor}.
\emph{seccomp} (secure computing mode) restricts which syscalls a process may
invoke~\cite{archwiki:seccomp}; a process that calls a blocked syscall is
terminated with \texttt{SIGSYS}. Together these three layers provide defense in
depth: an application that escapes its AppArmor profile still faces namespace
isolation; one that breaks out of the namespace still has a seccomp filter
limiting what kernel interfaces it can reach.

\begin{lstlisting}[language=bash, caption={Example permission file}]
# ~/.config/permissions/org.mozilla.firefox.toml

[filesystem]
allow = ["~/Downloads", "~/Documents", "~/.mozilla"]

[network]
allow = ["full"]

[devices]
allow = ["gpu"]

[graph]
# no entry = no access, never prompted
\end{lstlisting}

Graph access is never inferred or learned. An application with no
\texttt{[graph]} section has no graph access, with no exceptions and no
automatic prompting. Graph permissions must be explicitly declared.

\paragraph{Store-supplied profiles.}
The Store ships a curated permission profile alongside known applications.
Common software such as Firefox, LibreOffice, and VLC has community-maintained
profiles that are installed automatically and shown to the user at install time
for review. Community members can submit profiles for applications not yet
covered; submitted profiles go through Store review before publication, since
an overly permissive profile is a meaningful risk.

\paragraph{Flatpak adapter.}
Flatpak applications carry their own permission declarations. When a Flatpak is
installed, the system automatically generates a native permission file by
mapping Flatpak permissions to the native format: \texttt{filesystem=home}
becomes \texttt{filesystem.allow = ["\textasciitilde/"]}, and so on. The
generated file is shown to the user at install time and can be tightened before
confirming. The Flatpak manifest is the floor; the system never grants more
than what the Flatpak declared.

\paragraph{First-run learning mode.}
If no Store profile exists and the application is not a Flatpak, the first
launch runs in learning mode: all accesses are permitted and logged, nothing is
blocked. After the session ends, the user receives a summary of what the
application accessed, with options to review and confirm, edit first, or start
with all permissions blocked. From the second launch onward, the confirmed
profile is enforced. Learning mode applies to filesystem, network, and device
access only; graph access is always denied until explicitly granted and cannot
be learned.

\paragraph{Runtime behavior.}
When an application attempts an access outside its profile, the access is
denied and the user receives a notification offering to extend the profile.
If the user allows, the profile is updated and the application retries the
action; if the user denies, the block stands. Repeated denials for the same
path are collapsed: after three denials for the same access, the system asks
once whether to always block it rather than continuing to prompt.

\subsection{Module System}
\label{sec:modules}

The system is extensible at multiple layers through a unified module mechanism.
All module types use the same installation channel (the Store) and the same
manifest format, but they operate as independent subsystems with separate APIs.
A single Store package can bundle multiple module types together.

\paragraph{Module categories.}
Table~\ref{tab:modules} lists the available extension points. Theme modules
are a special case: they consist of CSS and token overrides only, contain no
executable code, and require no sandbox. All other categories run sandboxed
with declared permissions.

\begin{table}[htbp]
\caption{Module categories and extension points}
\label{tab:modules}
\begin{tabularx}{\linewidth}{llX}
\toprule
\textbf{Category} & \textbf{Extension point} & \textbf{Example} \\
\midrule
\texttt{theme}              & Design language layer    & Community color schemes \\
\texttt{spotlight.provider} & Spotlight search results & Web search, stock prices \\
\texttt{spotlight.action}   & Spotlight command actions & Create task, send message \\
\texttt{shell.topbar}       & Topbar element slot      & Custom status widget \\
\texttt{shell.widget}       & Desktop widget area      & Clock, weather, system stats \\
\texttt{mcp.server}         & AI MCP integration       & External service connector \\
\texttt{keybinding.profile} & Global input manager     & Alternative layout profiles \\
\bottomrule
\end{tabularx}
\end{table}

\paragraph{Manifest format.}
Every module package has a single manifest that declares what it provides,
its dependencies, conflicts, and permissions. The \texttt{[permissions]} block
uses the same schema as application permission files; at install time the
Module Runtime generates the permission file from the manifest and the user can
inspect and adjust it in Settings like any other application.

\begin{lstlisting}[language=bash, caption={Example module manifest}]
[package]
id      = "com.example.ai-spotlight"
name    = "AI Spotlight Provider"
version = "1.0.0"

[provides]
spotlight.provider = true

[dependencies]
shell   = ">=1.2"
modules = ["com.example.ai-core"]

[permissions]
network.allow = ["api.openai.com"]
graph.allow   = ["read"]
\end{lstlisting}

Modules cannot declare permissions beyond what the manifest specifies, and
manifests are reviewed during Store submission. Unlike applications, modules
cannot expand their permissions at runtime through a learning mode; their
profile is set at install time and any undeclared access returns
\texttt{EACCES}.

\paragraph{Module Runtime.}
A dedicated Module Runtime daemon manages lifecycle, sandboxing, and crash
recovery for all active modules. Each module runs in its own isolated WebView
context, separate from the shell's WebView. A module crash does not affect the
shell or other modules.

The runtime uses exponential backoff per module instance for crash recovery,
shown in Table~\ref{tab:crash-recovery}. The 60-second window resets on a
clean run, so a module that runs for several minutes and then crashes once
gets an immediate restart.

\begin{table}[htbp]
\caption{Module crash recovery policy}
\label{tab:crash-recovery}
\begin{tabularx}{\linewidth}{lX}
\toprule
\textbf{Crash count (within 60 s)} & \textbf{Behavior} \\
\midrule
1st & Immediate restart \\
2nd & Restart after 5 seconds \\
3rd & Restart after 30 seconds \\
4th or more & Marked as failed; no further automatic retries \\
\bottomrule
\end{tabularx}
\end{table}

When a module is marked failed, Settings shows a subtle error indicator with
the last crash log. The user can trigger a manual retry at any time. There is
no notification spam.

Module updates are downloaded in the background and applied at different points
depending on category: theme updates apply immediately as a token and CSS swap;
Spotlight modules update on the next Spotlight open; shell widgets and topbar
modules update on the next session start by default, with a manual
"Restart now" option available in Settings; MCP server modules update on the
next session start to avoid interrupting active AI interactions.

\paragraph{Shell API.}
Modules interact with the system through a restricted TypeScript SDK. No direct
access to the Smithay compositor or Graph Daemon is available; what is not in
the SDK is not accessible.

\begin{lstlisting}[language=bash, caption={Module SDK surface (TypeScript)}]
import { shell } from "@os/module-sdk"

// Window management
const windows = await shell.getWindows()
const active  = await shell.getActiveWindow()

// Graph (requires graph = "read" in manifest)
const recent = await shell.graph.query(
  "MATCH (f:File) RETURN f LIMIT 5"
)

// Spotlight provider registration
shell.spotlight.register({
  id:          "com.example.ai-spotlight",
  placeholder: "Ask AI...",
  onQuery:     async (q) => [ /* result objects */ ]
})
\end{lstlisting}

\paragraph{Global Input Manager.}
Keybindings and gestures are configured in a single place in Settings, not per
application or per module. Modules register named actions with default
bindings; the user overrides them centrally. The module never receives the
binding string directly; the Input Manager calls the registered action handler
when the binding fires. This means a user can rebind everything from one place
and no module can silently capture input outside its declared actions.

\paragraph{Developer experience.}
A CLI scaffolds a complete module project with one command:
\texttt{npx create-os-module <n>}, producing a \emph{Vite}~\cite{vite},
TypeScript, Tailwind, and module-sdk setup with a manifest template.
Vite is a front-end build tool and development server; its key property is a
near-instant dev server that uses the browser's native ES module support to
serve individual files without bundling, making hot-reload on code changes
nearly instantaneous regardless of project size.

In dev mode, the module runs unsigned in a dedicated dev sandbox, all declared
permissions are granted without Store review, WebView DevTools are accessible,
and hot-reload via Vite is active. Dev mode is only available on developer
builds; production builds reject unsigned modules.

For distributing a finished module outside the Store, a developer generates a
self-signed key pair with the CLI, signs the module, and distributes it as an
\texttt{.osmod} file. Installation uses \emph{TOFU} (Trust On First Use): no
certificate authority pre-approves the key; the first install prompts the user
to trust the key fingerprint explicitly, and all subsequent updates signed with
the same key are accepted silently. This is the same model used by SSH host
key verification. Sideloaded modules receive no expanded permissions and are
marked ``Not from Store'' in Settings.

\subsection{Store}
\label{sec:store}

The Store is the primary distribution channel for applications, themes,
modules, and profile packages. It is designed to be financially sustainable for
the project without becoming a gatekeeper for the ecosystem.

\paragraph{Publishing.}
Publishing an application to the Store is free and remains free indefinitely.
There is no annual developer fee. Every upload runs automated technical checks:
valid permission declarations, correct package signing, sandbox policy
compatibility, and resolvable declared dependencies. These checks are
automated; no manual review is required for packages that pass them. A failed
check returns a specific error. A manual Security Review is available
optionally and is the basis for the Security Audit Badge described below.

\paragraph{Revenue model.}
Supported pricing models include free, one-time purchase, subscription,
in-app purchases, pay-what-you-want, and freemium. Developers can also direct
users to external payment links; the Store payment infrastructure is a
convenience, not a monopoly.

Commission is calculated on total annual revenue across all applications in a
developer account, using the tiered structure in Table~\ref{tab:commission}.
When a threshold is crossed mid-year, the new rate applies to all subsequent
transactions in that year, not retroactively.

\begin{table}[htbp]
\caption{Store commission tiers}
\label{tab:commission}
\begin{tabularx}{\linewidth}{XX}
\toprule
\textbf{Annual revenue (account total)} & \textbf{Commission} \\
\midrule
Up to \EUR{10,000}                        & 0\% \\
\EUR{10,000} to \EUR{50,000}              & 5\% \\
\EUR{50,000} to \EUR{500,000}             & 15\% \\
\EUR{500,000} to \EUR{5,000,000}          & 20\% \\
Over \EUR{5,000,000}                      & 25\% \\
\bottomrule
\end{tabularx}
\end{table}

\paragraph{User-facing features.}
Purchases are tied to the user account, not the device. Trials are declared in
the app manifest and enforced by the Store; no credit card is required to start
a trial, and when it ends the user receives a single notification with a
purchase option and no automatic charges. Family Sharing is opt-in per
application: developers can allow up to six to eight family members to use an
application under one purchase. Every application's Store page shows a full
price history including every price change with its date; this is automatic and
cannot be disabled by the developer. All active subscriptions across all
applications are visible and cancellable in one place in Settings.

\paragraph{Developer dashboard.}
Developers see install counts, active user numbers, revenue and payout history,
ratings and written reviews, crash reports from users who opted into crash
reporting at the system level, and version adoption numbers. All data is
aggregated and anonymized; the developer sees counts, never identities. No
third-party analytics are required or supported.

\paragraph{Security Audit Badge.}
Applications can display a Security Audit badge on their Store page by
uploading an independent security audit report. The Store verifies that the
report originates from the stated auditing organization and links it from the
application page. Users can read the full report and judge the auditing
organization's credibility themselves. The badge means an audit exists and is
accessible, not that the Store has endorsed the application's security.

\paragraph{Long-term support.}
Developers can declare a Long Term Support commitment in the application
manifest by setting a support end date. The Store displays this date publicly
on the application page. It is a voluntary public commitment with no automated
enforcement, but the public visibility creates accountability. LTS applications
are filterable in Store search for users who prioritize stability.

\paragraph{Profile package distribution.}
Signed environment profile packages (\texttt{.lenv} files, described in
Section~\ref{sec:platform}) are distributed through the Store using the same
infrastructure as applications. Profile packages are free to publish and carry
no commission.
